\documentclass[
    12pt,
    a4paper
]{article}

% ==================================================
% Idioma e encoding
% ==================================================
\usepackage[T1]{fontenc}
\usepackage[utf8]{inputenc}
\usepackage{lmodern}
\usepackage[portuguese]{babel}
\usepackage{csquotes}

% ==================================================
% Página
% ==================================================
\usepackage[
    top=2.5cm,
    bottom=2.5cm,
    left=2.5cm,
    right=2.5cm
]{geometry}
\usepackage{fancyhdr}
\usepackage{lastpage}

% ==================================================
% Texto
% ==================================================
\usepackage{setspace}
\setstretch{1.15}

\usepackage{microtype}

% ==================================================
% Matemática e unidades
% ==================================================
\usepackage{amsmath}
\usepackage{amssymb}
\usepackage{siunitx}

% ==================================================
% Figuras e tabelas
% ==================================================
\usepackage{graphicx}
\usepackage{booktabs}
\usepackage{float}
\usepackage{tabularx}

\graphicspath{{figures/}}

% ==================================================
% Código
% ==================================================
\usepackage{xcolor}
\usepackage{listings}

\lstset{
    basicstyle=\ttfamily\small,
    frame=single,
    breaklines=true,
    numbers=left,
    numberstyle=\tiny,
    showstringspaces=false
}

% ==================================================
% Links
% ==================================================
\usepackage[
    colorlinks=true,
    linkcolor=blue,
    citecolor=blue,
    urlcolor=blue
]{hyperref}

% ==================================================
% Bibliografia
% ==================================================
\usepackage[
    backend=biber,
    style=ieee,
    sorting=none
]{biblatex}

\addbibresource{references/referencias.bib}

% ==================================================
% Documento
% ==================================================

% Dados do relatório: editar estes comandos para atualizar a capa.
\newcommand{\tituloRelatorio}{Sistema de Escrita/Leitura Multi-Protocolo}
\newcommand{\subtituloRelatorio}{I2C, RS232, RS485, CAN e Ethernet}
\newcommand{\disciplinaRelatorio}{LASET --- Laboratório de Sistemas Autónomos}
\newcommand{\cursoRelatorio}{Nome do Curso}
\newcommand{\autoresRelatorio}{Francisco Liquito\\Pedro Ribeiro}
\newcommand{\anoLetivoRelatorio}{2026/2027}

\definecolor{corCapa}{HTML}{17324D}

\pdfstringdefDisableCommands{\def\\{; }}
\hypersetup{
    pdftitle={\tituloRelatorio},
    pdfauthor={\autoresRelatorio}
}

% Paginação do relatório; a capa não apresenta cabeçalho nem rodapé.
\pagestyle{fancy}
\fancyhf{}
\renewcommand{\headrulewidth}{0pt}
\renewcommand{\footrulewidth}{0pt}
\fancyfoot[C]{\small Página \thepage\ de \pageref*{LastPage}}
\fancypagestyle{plain}{%
    \fancyhf{}
    \renewcommand{\headrulewidth}{0pt}
    \renewcommand{\footrulewidth}{0pt}
    \fancyfoot[C]{\small Página \thepage\ de \pageref*{LastPage}}
}

% ==================================================
% Início do documento
% ==================================================

\begin{document}

% --------------------------------------------------
% Capa
% --------------------------------------------------

\hypersetup{pageanchor=false}
\begin{titlepage}
    \thispagestyle{empty}
    \centering
    \sffamily

    {\large\bfseries \cursoRelatorio\par}
    \vspace{0.4cm}
    {\large \disciplinaRelatorio\par}

    \vfill

    {\small\bfseries\color{corCapa} RELATÓRIO TÉCNICO\par}
    \vspace{0.6cm}
    {\color{corCapa}\rule{\textwidth}{1.2pt}\par}
    \vspace{0.8cm}
    {\Huge\bfseries\color{corCapa} \tituloRelatorio\par}
    \vspace{0.6cm}
    {\Large \subtituloRelatorio\par}
    \vspace{0.8cm}
    {\color{corCapa}\rule{\textwidth}{0.4pt}\par}

    \vfill

    {\small\bfseries AUTORES\par}
    \vspace{0.4cm}
    {\large \autoresRelatorio\par}

    \vspace{1.8cm}
    {\small Ano Letivo\par}
    \vspace{0.2cm}
    {\large\bfseries \anoLetivoRelatorio\par}
\end{titlepage}
\hypersetup{pageanchor=true}

% --------------------------------------------------
% Índice
% --------------------------------------------------

\tableofcontents

\newpage

% ==================================================
% Secções
% ==================================================

\section{Introdução}
Este documento descreve o desenvolvimento de um sistema de
escrita e leitura multiprotocolo com suporte para I2C, RS232,
RS485, CAN e Ethernet.


\section{Problema e Motivação}
Os sistemas autónomos e os sistemas embebidos integram funções de
perceção, processamento, controlo e atuação que dependem da troca de
informação entre os seus elementos. A aquisição de dados de sensores,
o envio de comandos a atuadores e a comunicação com controladores ou
equipamentos de supervisão podem envolver tecnologias de comunicação
distintas numa mesma instalação. Esta coexistência resulta das
características dos dispositivos e das funções que desempenham, bem
como da necessidade de integrar equipamentos desenvolvidos em momentos
ou contextos diferentes.

Neste contexto, sensores, atuadores, controladores e outros equipamentos
podem disponibilizar interfaces como I2C, RS232, RS485, CAN e Ethernet.
Estas tecnologias apresentam diferenças nas características elétricas,
na organização das ligações e nos mecanismos de comunicação que
suportam. Por conseguinte, a ligação de um novo equipamento exige a
análise da compatibilidade da interface e da forma como os dados são
transmitidos e recebidos, além da compreensão do seu significado para
a aplicação.

O problema central reside na heterogeneidade destas interfaces e na
dificuldade de integrar operações de leitura e escrita sobre equipamentos
com formas de acesso distintas. O suporte de cada tecnologia pode exigir
configurações, procedimentos de comunicação e mecanismos de tratamento
de erros específicos. Quando estas particularidades são tratadas de
forma dispersa na aplicação, a inclusão ou substituição de equipamentos
pode implicar alterações em várias partes do sistema, dificultando a
manutenção e a reutilização das funções de comunicação. Além disso, a
disponibilidade de uma ligação compatível não assegura, por si só, a
interpretação correta dos dados trocados.

A relevância do trabalho reside, assim, na possibilidade de organizar de
forma coerente os mecanismos de acesso aos equipamentos, facilitando a
integração e a evolução de sistemas com diferentes tecnologias de
comunicação.


\section{Estado da Arte}
\subsection{I2C}

O protocolo I2C é frequentemente utilizado para comunicação entre
circuitos integrados e periféricos em sistemas embebidos.

\subsection{RS232}

A interface RS232 é tradicionalmente utilizada em comunicações
série ponto-a-ponto.

\subsection{RS485}

A interface RS485 utiliza transmissão diferencial e é amplamente
utilizada em ambientes industriais.

\subsection{CAN}

O barramento CAN permite a comunicação entre múltiplos nós e é
utilizado em aplicações automóveis e industriais.

\subsection{Ethernet}

Ethernet permite a comunicação de dispositivos através de redes
locais, suportando protocolos de níveis superiores como IP, TCP e
UDP.

\section{Requisitos}
\subsection{Requisitos Funcionais}

Os requisitos funcionais seguintes decorrem do título do trabalho e
delimitam as capacidades de comunicação a desenvolver.

\begin{description}
    \item[RF01 --- Suporte para I2C:]
    O sistema deve suportar comunicação através de uma interface I2C.

    \item[RF02 --- Suporte para RS232:]
    O sistema deve suportar comunicação através de uma interface RS232.

    \item[RF03 --- Suporte para RS485:]
    O sistema deve suportar comunicação através de uma interface RS485.

    \item[RF04 --- Suporte para CAN:]
    O sistema deve suportar comunicação através de uma interface CAN.

    \item[RF05 --- Suporte para Ethernet:]
    O sistema deve suportar comunicação através de uma interface Ethernet.

    \item[RF06 --- Escrita e transmissão:]
    O sistema deve permitir operações de escrita ou transmissão de dados
    através de cada uma das cinco interfaces previstas, de acordo com as
    características da respetiva tecnologia.

    \item[RF07 --- Leitura e receção:]
    O sistema deve permitir operações de leitura ou receção de dados
    através de cada uma das cinco interfaces previstas, de acordo com as
    características da respetiva tecnologia.
\end{description}

\subsection{Requisitos Não Funcionais}

Os requisitos não funcionais seguintes constituem propostas preliminares,
a validar com os docentes. Os respetivos critérios de aceitação serão
definidos após a consolidação dos requisitos.

\begin{description}
    \item[RNF01 --- Modularidade:]
    A organização do sistema deverá permitir separar as funções associadas
    às diferentes interfaces de comunicação.

    \item[RNF02 --- Facilidade de manutenção:]
    A implementação deverá ser organizada e documentada de forma a
    facilitar a compreensão, a correção e a alteração das suas funções.

    \item[RNF03 --- Extensibilidade:]
    A arquitetura deverá facilitar a introdução futura de novas
    funcionalidades de comunicação, sem exigir a reformulação integral
    do sistema.

    \item[RNF04 --- Fiabilidade:]
    O sistema deverá executar as operações de comunicação de forma
    consistente nas condições de funcionamento que vierem a ser definidas.

    \item[RNF05 --- Facilidade de teste e diagnóstico:]
    A organização do sistema deverá facilitar a verificação das operações
    de comunicação e a identificação de falhas.
\end{description}

\subsection{Restrições do Projeto}

O âmbito conhecido do projeto abrange operações de leitura e escrita
através de I2C, RS232, RS485, CAN e Ethernet. Não estão ainda confirmadas
restrições técnicas específicas relativas à implementação.

A plataforma de execução, a organização concreta das interfaces e as
eventuais funções de interação entre tecnologias permanecem por definir.
Esta ausência de definição corresponde ao estado preliminar do projeto
e não constitui, por si só, uma restrição técnica. As restrições
adicionais deverão ser registadas após confirmação com os docentes.

\subsection{Requisitos a Validar com os Docentes}

As questões seguintes destinam-se a esclarecer o âmbito do sistema e
a estabelecer requisitos verificáveis, sem antecipar decisões de
implementação.

\begin{description}
    \item[QV01 --- Interação entre tecnologias:]
    É necessária conversão ou encaminhamento de dados entre protocolos,
    ou apenas a realização de operações de leitura e escrita através
    de cada interface?

    \item[QV02 --- Protocolos normalizados:]
    É obrigatória a utilização de protocolos normalizados, como Modbus?
    Em caso afirmativo, quais?

    \item[QV03 --- Plataforma de execução:]
    Existe um requisito de utilização de microcontrolador ou é permitida
    uma plataforma com Linux?

    \item[QV04 --- Variante de CAN:]
    O suporte para CAN deve abranger Classical CAN, CAN FD ou ambas
    as variantes?

    \item[QV05 --- Âmbito de Ethernet:]
    O que se entende por suporte para Ethernet neste projeto:
    apenas a camada física ou também comunicação através de TCP,
    UDP ou protocolos de aplicação?

    \item[QV06 --- Placa de circuito impresso:]
    É necessário desenvolver uma placa de circuito impresso
    (PCB) própria?

    \item[QV07 --- Desempenho:]
    Existem métricas obrigatórias de desempenho? Em caso afirmativo,
    quais são os respetivos critérios de aceitação e condições
    de avaliação?
\end{description}


\section{Arquitetura de Alto Nível}
Nesta fase será definida apenas a arquitetura de alto nível do
sistema, sem selecionar componentes de hardware específicos.

% ==================================================
% Bibliografia
% ==================================================

\newpage

\printbibliography

\end{document}